\subsection{The column index away from the prime assigned to the sequence}
\label{subsec:column-orbit-length}

Fix the word and its level-zero class model. At a coefficient place over \(p\ne T_u\), the neighbor tests on the selected Frobenius orbits determine a bound for the length on the trace-zero column.
The set \(I\) of selected Frobenius orbits, their orders \(t_v\), the
common test depth \(D\), and their deficient fractions \(\kappa_v\)
are those of Sections~\ref{subsec:frobenius-selection}
and~\ref{sec:host-edge-tests}. Here \(t_v\) denotes the nonnegative
integer order cutoff, rounded up at each orbit. On the primes with selected orbits
the weight selection gives \(\sigma=m\sum_{v\in I}\kappa_v\le1/2\).

For each \(v\) the nonnegative integer \(S_v\) is the ceiling of the
sum over the \(m\) stages of the additional gain at \(v\) obtained
by using that orbit. More explicitly, if \(g^U_{jv}\) is the up
gain, \(P_{jv}\) its price and \(g^0_{jv}\) the base gain, then
\[
 S_v=\left\lceil\sum_{j=1}^m
       \big(g^U_{jv}-P_{jv}-\max(0,g^0_{jv})\big)_+\right\rceil .
\]
Write \(S_*=\max_{v\in I}S_v\), with value zero on the empty set.

Temporarily work over the common local coefficient ring \(O\) of
Section~\ref{sec:host-edge-tests}; write \(\omega^2=p\),
\(c_p=\operatorname{length}_O(O/pO)\).
At a vertex in one orbit the three value lattices are
\[
 A_v=N_-,\qquad H_v=H(L),\qquad C_v=(N_+)^\vee.
\]
The dual in the last expression is taken in the common polynomial
space with its evaluation pairing. They satisfy
\[
 C_v\subseteq H_v\subseteq A_v,\qquad
 \omega^{t_v}A_v\subseteq H_v,\qquad
 \omega^{t_v}H_v\subseteq C_v .
 \tag{C1}\label{eq:tag-c1}
\]
\begin{proposition}[Neighbor images in the orbit value lattices]\label{prop:orbit-value-lattice-neighbors}
With the selected level-zero orbit lattices, the normalized
incoming sums preserve \(M_0\), send \(A_v\) into \(C_v\) at
the chosen factor and commute at distinct \(v\). The actual
lattice \(M\) has the comparison \textup{\eqref{eq:tag-c2}}. For every
nonempty \(B\subseteq I\) there is a fixed exceptional set in
\(M_B\) satisfying \textup{\eqref{eq:tag-c3}}, and its cardinality satisfies
\[
 \#\mathcal C_B\le d\prod_{v\in B}\kappa_v.
\]

\end{proposition}

\begin{proof}
\proofstep{Use the diagonal models for the edge and its dual.}

Indeed \(N_-\) lies in \(\omega^{-t_v}H_v\) and contains \(H_v\).
The dual statements follow from

\[
  H_v^\vee\subseteq N_+\subseteq\omega^{-t_v}H_v^\vee.
\]
For the edge constraints in these containments choose the diagonal
basis adapted to a neighbor. In a component the one-parameter
diagonal weight on its first \(r\) entries is at most the sum of
the first \(r\) entries of the dominant weight, by the weight
inequalities; on the last lattice take the dual inequalities in
the inverse transpose frame.
Thus for a torus weight \(\nu\) occurring in the component of highest weight \(\lambda\), the bound on the selected diagonal direction is
\[
  \sum_{i=1}^r\nu_i\le\sum_{i=1}^r\lambda_i .
\]
 Summing over the Frobenius components
is exactly the normalizing exponent \(A_r\) in \(R_e,R'_e\).
Thus the two incident ordinary lattices lie in the indicated
normalized edge and dual edge. These inclusions include the
determinant powers for signed weights.

\proofstep{Normalize each incoming sum on the class values.}

On the lattice of functions on the level classes let \(Y_v\) be
\(p^{A_r}\) times the incoming-neighbor sum at \(v\). The unnormalized
sum leaves every other vertex and coefficient factor fixed.
The operators commute, since the two finite choices at distinct
places are independent.
For distinct places the double sum is over a Cartesian product, hence
\[
  Y_vY_{v'}=Y_{v'}Y_v\qquad(v\ne v').
\]

For any mixture of value lattices at other
places, \(Y_v\) sends the choice \(A_v\) into \(C_v\). Indeed a source
value belongs to \(R_e\) and the target test belongs to \(R'_e\);
their pairing times \(p^{A_r}\) is integral on every edge.
It also preserves the all-ordinary lattice: in the same diagonal
bases the map from source to target involves powers with denominator
at most \(p^{A_r}\). On the divided-power Weyl sublattice this
follows in a weight basis, since the basis changes are integral
automorphisms and the torus acts by its weights.

\proofstep{Compare the standard and divided-power Weyl lattices.}

Let \(M_0\) be the class lattice with the value lattices \(H_v\) at
all places and let \(M\) be the actual level-zero Weyl lattice.
At sufficiently large \(x\) its defining integer \(e_0\) satisfies
\[
 p^{e_0}M_0\subseteq M\subseteq M_0,\qquad
      e_0\ll f k^2\log_p a\quad (p\le a).
 \tag{C2}\label{eq:tag-c2}
\]
For a spread \(l\) on one embedding the trace comparison in
Section~\ref{sec:edge-contraction} puts \(p^{\binom k2
\lfloor\log_p(l+k-1)\rfloor}\) times the Schur lattice in the
divided-power span, which lies in the Schur lattice. Tensor these
inclusions over the \(f\) positions. The host bounds give

\[
  l+k-1\le a^2,
\]
 and \(e_0\) is the chosen exponent dominating
this sum. Taking a ceiling with fixed coefficient changes the bound
by at most one, absorbed since \(\log_p a\ge1\). Conjugating a
frame acts on both lattices; the comparison therefore holds for all
class values. Both lattices have rank \(d=d_S\).

\proofstep{Tensor the retained Smith bases at the selected vertices.}

For \(B\subseteq I\), put \(M_B\) equal to the class lattice with
values \(C_v\) at \(v\in B\) and \(H_v\) elsewhere. Write
\(M_{B,v}^{+}\) for the same choice with \(C_v\) replaced by
\(A_v\) at \(v\in B\). For each nonempty \(B\) there is a fixed
set \(\mathcal C_B\subset M_B\) with
\[
 \begin{split}
 |\mathcal C_B|&\le d\prod_{v\in B}\kappa_v,\\
 p^{D+1}M_B&\subset
       \sum_{v\in B}Y_v M_{B,v}^{+}+O\mathcal C_B .
 \end{split}                                                     \tag{C3}\label{eq:tag-c3}
\]
Here and below the \(O\)-multiple of a set denotes its linear span.

To prove this, take orbit representatives for the finite class set.
Their arithmetic stabilizers are trivial by the type level argument
at the beginning of Proposition~\ref{prop:trace-zero-column}. An elementary
value at a representative defines a section supported on its orbit:
extend it by the representation action and the unit level character.
The extension is well defined by freeness, recovers the value at the
representative, and gives independent sections from distinct classes.
In particular, if \(s_a\) is supported on the orbit of representative \(a\), the evaluation at representatives is diagonal:
\[
  s_a(b)=\begin{cases}s_a(a)&b=a,\\0&b\ne a.\end{cases}
\]

At a vertex, take the Smith basis of the target \(C_v\) given by
the rank computation in Section~\ref{sec:host-edge-tests}.
In more detail, the normalized incoming columns \(p^{A_r}F_e\)
are integral in \(C_v\), and their evaluations against elements of
\(N_+\) contain a \(p^D\)-scaled identity of the retained size.
A Smith form for the columns in \(C_v\) therefore has at least
that many elementary exponents at most \(v_O(p^D)\).
Indeed a map with this scaled right inverse surjects onto that
many coordinates after allowing the factor \(p^D\); elementary
divisors over the DVR give the assertion.

For any retained Smith coordinate of exponent \(e\), the required divisibility is
\[
  e\le v_O(p^D)\quad\Longrightarrow\quad p^DO\subseteq\mathfrak m_O^e .
\]
 Multiplication by
\(p^D\) on each such basis vector is consequently an integral
combination of these incoming columns. The remaining indices form
at most the fraction \(\kappa_v\) of the basis. At other factors
use full ordinary bases.

Tensor the local bases at each class representative \(c\). Write
\(d_c\) for its full tensor rank and let \(\mathcal C_{B,c}\)
contain the products with a remaining Smith index at every
\(v\in B\). Independence of the factors gives
\[
 |\mathcal C_B|=\sum_c|\mathcal C_{B,c}|
          \le \sum_c d_c\prod_{v\in B}\kappa_v
          =d\prod_{v\in B}\kappa_v .
\]
Every other product has a retained factor at some \(v\in B\).
Put its incoming source values in orbit sections with all other tensor factors
unchanged, and apply \(Y_v\).

At the target this gives the normalized
column. Every other contribution is a transport from a competing
edge and is divisible by \(p^{D+1}\) in the target lattice.
Indeed the source test gives order \(c_*=8(D+1)\) in the competing
normalized edge lattice. Pair it against \(N_+\) to get the same
order in the predual, using 
\[
  \omega^2=p.
\]
This also treats distinct incoming terms landing in the same class.
Only a transport carrying the source and its distinguished target
to the prescribed pair can supply the distinguished term. Two such
transports would differ by a stabilizer of the target. The other
terms each use a competing edge after transport back to the source.

\proofstep{Remove the higher-order residual terms.}

Let \(J_B=\sum_{v\in B}Y_vM_{B,v}^{+}+O\mathcal C_B\).
We have shown 
\[
  p^D M_B\subset J_B+p^{D+1}M_B.
\]
The image of \(p^D M_B\) in \(M_B/J_B\) is a finite \(O\)-module
in the sense of finite generation, since there are finitely many
class representatives and all value lattices are finite free.
Put \(D_B=(p^DM_B+J_B)/J_B\). The inclusion above gives
\[
 D_B\subseteq(p^{D+1}M_B+J_B)/J_B=pD_B\subseteq D_B .
\]
Nakayama's lemma therefore gives \(D_B=0\), giving
\eqref{eq:tag-c3}, including the stated allowance of one additional power.

\end{proof}

\subsubsection{Large elementary divisors}

A bound in the full rational span must survive passage to the trace-zero lattice. This requires estimating the elementary divisors beyond the selected depth.

For either lattice \(L=M_0,M\) define
\[
 P(L)=\sum_{J\subseteq I}
       p^{\sum_{v\notin J} S_v}
              \left(\prod_{v\in J}Y_v^m\right)L,\qquad
 \tau_0=m\sum_{v\in I}(t_v+D+1).
\]
\begin{proposition}[Tail count for the elementary divisors]\label{prop:orbit-smith-tail}
For each integer \(\beta\ge0\), \(p^{\tau_0+\beta S_*}M_0\)
is contained in \(P(M_0)\) plus the span of the number of fixed
elements in \textup{\eqref{eq:tag-c4}}. If \(I\ne\varnothing\) and
\(\tau=e_0+\tau_0\), this also bounds the number of elementary
exponents of \(P(M)\subseteq M\) strictly above
\(c_p(\tau+\beta S_*)\). Writing \(\alpha_i\) for these elementary exponents with multiplicities, one has
\[
 \#\{i:\alpha_i>c_p(\tau+\beta S_*)\}\le \frac{d\,\sigma^{\beta+1}}{(\beta+1)!}.
\]

\end{proposition}

\begin{proof}
The power notation on \(Y_v\) means iteration. Write \(S_*=0\) if
\(I\) is empty. We prove that for every \(\beta\ge0\),
\(p^{\tau_0+\beta S_*}M_0\) lies in \(P(M_0)\) plus the span of
at most
\[
             d\,m^{\beta+1}
               \sum_{\substack{B\subset I\\|B|=\beta+1}}
                     \prod_{v\in B}\kappa_v                      \tag{C4}\label{eq:tag-c4}
\]
fixed elements of \(M_0\). In this sum \(B=I\) is allowed.

\proofstep{Iterate with the first unfinished coordinates.}

Order \(I\), and attach to a term the exponents \(0\le a_v\le m\)
of the operators \(Y_v\) already applied outside its argument.
A coordinate is finished when \(a_v=m\). If there are at least
\(\beta+1\) unfinished coordinates, use their first \(\beta+1\)
as \(B\) and apply Proposition~\ref{prop:orbit-value-lattice-neighbors}
to the argument using \textup{\eqref{eq:tag-c3}}. In the summand indexed
by \(v\), one more \(Y_v\) moves outside the argument and \(a_v\)
increases by one.
Thus the exponent vector attached to that summand changes by
\[
  (a_i)_{i\in I}\longmapsto(a_i+\mathbf1_{\{i=v\}})_{i\in I}.
\]
 The new argument is at \(A_v\); it is at \(C\)
at all the other coordinates of \(B\).

If \(v\) is unfinished
and the procedure continues, multiply by \(\omega^{2t_v}\) to
return that coordinate to \(C_v\). If \(v\) is finished, return it
to \(H_v\) by \(\omega^{t_v}\) and introduce the next unfinished
coordinate, if needed, by the move \(H_i\to C_i\) costing
\(\omega^{t_i}\). Use the latter move also at the first entry of
each coordinate. Other coordinates of \(B\) keep their lattices.

\proofstep{Bound the cost along each branch.}

Work along a branch with the fixed initial scalar
\(p^{\tau_0+\beta S_*}\), distributing the unused part to
every summand. Each coordinate enters \(B\) at most once.
Its transport costs, in powers of \(\omega\), are bounded by
\[
 \begin{cases}
 t_v+2(m-1)t_v+t_v=2mt_v,&\text{after all \(m\) increments},\\
 (1+2j)t_v\le2mt_v,&\text{if stopped after \(j<m\) increments}.
 \end{cases}
\]
The second bound includes any final return from \(A_v\) to
\(H_v\); the other open lattices \(C_i\) are already in \(H_i\).
There are at most \(m|I|\) applications of \eqref{eq:tag-c3}: an exceptional
term at the last application uses up a possible increment.
The budget for increments comes from
\[
  0\le a_i\le m,\qquad \sum_{i\in I}a_i\le m|I| .
\]
Thus \(p^{\tau_0}\) pays all costs from \eqref{eq:tag-c1} and \eqref{eq:tag-c3}.

\proofstep{Finish inside the product module.}

Once at most \(\beta\) coordinates remain unfinished, use the
scalar \(p^{S_i}\) at each unfinished coordinate. The product
costs at most \(p^{\beta S_*}\) in the sense of divisibility.
Let \(F\subset I\) be the finished coordinates. Commutativity separates the completed factors from the operators at the unfinished coordinates
\[
  \prod_{i\in I}Y_i^{a_i}=\left(\prod_{i\in F}Y_i^m\right)\left(\prod_{i\in I\setminus F}Y_i^{a_i}\right),\qquad \left(\prod_{i\in I\setminus F}Y_i^{a_i}\right)M_0\subseteq M_0.
\]
 The resulting term belongs
to \(P(M_0)\). This covers the initial case \(|I|\le\beta\).

\proofstep{Count the exceptional rank and descend to \(M\).}

It remains to count exceptional terms. For a fixed active set \(B\)
of size \(\beta+1\), every coordinate before its last member and
outside \(B\) is finished. Every coordinate after its last member
has exponent zero, since it could never have belonged to an earlier
active set. The outside factors are thus determined by \(B\).
Inside \(B\) there are just \(m^{\beta+1}\) possible exponent lists.
Equivalently the number of exponent choices on the active set is
\[
  \prod_{v\in B}m=m^{|B|}=m^{\beta+1}.
\]
The exceptional generators are those products of commuting \(Y_i\)
applied to \(\mathcal C_B\); any unused scalar multiplies them
integrally. They belong to \(M_0\). Counting them proves \eqref{eq:tag-c4}.
This argument needs simultaneous tests at the different places:
the product of deficiencies comes from the tensor values in \eqref{eq:tag-c3}.

Let \(I\ne\varnothing\) and put \(\tau=e_0+\tau_0\). By \eqref{eq:tag-c2},

\[
  p^{e_0}P(M_0)\subset P(M);
\]
 multiplying all the exceptional
vectors by \(p^{e_0}\) puts them in \(M\). Consequently \eqref{eq:tag-c4}
also bounds the number of elementary divisors of \(P(M)\) in
\(M\) with exponent larger than \(c_p(\tau+\beta S_*)\).
In fact \(p^{\tau+\beta S_*}(M/P(M))\) lies in a submodule with
at most the displayed number of generators; a submodule over a DVR
of a module with that many generators also has at most that many.
The empty subset in \(P(M)\) gives \(p^{\sum S_v}M\subset P(M)\),
so all exponents are finite. Each squarefree monomial appears \(s!\) times in the power sum
expansion. Since all \(\kappa_v\ge0\),
\[
 s!\sum_{\substack{B\subset I\\|B|=s}}\prod_{v\in B}\kappa_v
              \le \left(\sum_{v\in I}\kappa_v\right)^s
                  \qquad(s\ge1).
\]
Taking \(s=\beta+1\) bounds \eqref{eq:tag-c4} by
\(d\,\sigma^{\beta+1}/(\beta+1)!\).

\end{proof}

Write \(V_{\rm ram}\) for the local trace-zero vector space and
\(M_{\rm ram}=M\cap V_{\rm ram}\). The local unshifted column is
\[
 N_w=P(M)\cap p^\tau M\cap V_{\rm ram}.
\]
Saturation of the intersection with a vector space gives
\(N_w\subset p^\tau M_{\rm ram}\). Moreover \(M_{\rm ram}/N_w\)
injects into \(M/(P(M)\cap p^\tau M)\).

\begin{lemma}[Length on the local trace-zero intersection]\label{lem:orbit-trace-zero-colength}
Set \(M_{\rm ram}=M\cap V_{\rm ram}\) and
\(N_w=P(M)\cap p^\tau M\cap V_{\rm ram}\) for the local trace-zero
space \(V_{\rm ram}\). The elementary divisor bound of
Proposition~\ref{prop:orbit-smith-tail} gives
\(N_w\subset p^\tau M_{\rm ram}\) and
\[
 \operatorname{length}_O(M_{\rm ram}/N_w)/c_p
   \le d(\tau+2 S_*\sigma).
\]
When the set of orbits is empty, \(N_w=M_{\rm ram}\) and
the length is zero.
\end{lemma}

\begin{proof}
In a Smith basis let
\(d_i\) denote the elementary exponents of \(P(M)\) in \(M\).
The latter length is \(\sum_i\max(d_i,c_p\tau)\).
Count its excess over \(dc_p\tau\) in intervals of width \(c_p S_*\)
using \eqref{eq:tag-c4}. When \(S_*=0\), the empty subset already gives \(P(M)=M\).
For \(S_*>0\) the sum is bounded by
\(dc_pS_*\sum_{h\ge1}\sigma^h/h!\le 2dc_p S_*\sigma\), since
\(0\le\sigma\le1/2\). Hence
\[
  N_w\subset p^\tau M_{\rm ram},\qquad
  \frac{\operatorname{length}_O(M_{\rm ram}/N_w)}{c_p}
        \le d(\tau+2S_*\sigma).                                  \tag{C5}\label{eq:tag-c5}
\]
For \(I=\varnothing\) the selected convention is \(\tau=0\),
\(P(M)=M\), \(N_w=M_{\rm ram}\); \eqref{eq:tag-c5} then holds with equality
in its length bound. In particular the length vanishes at every
place without the local orbit condition away from \(T_u\).
\end{proof}

\subsubsection{The local index on the selected lattice}

The preceding tail estimate can now be applied to the actual level-zero lattice. The level trace and the prescribed orbit cuts determine which local index enters the column slope.

\begin{proposition}[Normalized index and excess in the original column]\label{prop:orbit-local-descent}
The local trace-zero lattice in Lemma~\ref{lem:orbit-trace-zero-colength}
descends with the same normalized colength to the actual
coefficient place. On the selected family its excess cost
satisfies \textup{\eqref{eq:tag-c6}} on every subset of served rational primes,
and beyond the family threshold the baseline satisfies
\[
 2\sum_p\tau(p)\log p\le b^2.
\]

\end{proposition}

\begin{proof}
\proofstep{Descend the operators, intersections and normalized lengths.}

Apply Lemma~\ref{lem:orbit-trace-zero-colength} with the
trace-zero intersection obtained after base change.
The local coefficient extension does not change the normalized
length in \eqref{eq:tag-c5}. Here is why the module being compared descends.
The level-zero basis maps identify the rational class values with
polynomials over the coefficient field \(K\), evaluated in the
chosen local frames via their coefficient embedding. On a valid tuple
the identification follows by writing the algebraic section on the
translated rational representative, with its rational level character;
off the set of valid tuples it is zero. One may recover the rational
subspace inside its extension from this test on all coefficients:
choose a nonzero square minor of the array of coefficient evaluations
of the basis functions.

Such a minor exists by their independent
images. With \(A\) the chosen invertible coordinate minor of size \(n\) and \(y\) the restricted vector of rational values, the coefficient vector is
\[
  c=A^{-1}y\in K^n,
\]
since both \(A\) and \(y\) have entries in \(K\).
The neighbor sum has finitely many terms and scalar \(p^{A_r}\in K\)
at every tuple. It preserves the generic span, since it preserves
the integral value lattice, and the coefficient test shows that its
action descends as a \(K\)-linear operator. The trace condition
likewise comes from the rational diamond operators. The local model
of \(M\) is the \(O\)-span of the level-zero lattice in these
coordinates. This descent uses the class-value identification for the
\emph{same} word and level as the trace condition.

It follows that \(P(M)\) is obtained by base change: each summand
in its definition is the image of an integral scalar times a
descending rational operator on the level-zero lattice.
Flat scalar extension preserves each image in the sum, by the exact identity
\[
  \operatorname{im}(A)\otimes O=\operatorname{im}(A\otimes O).
\]
Intersection with \(p^\tau M\) and with the rational trace-zero
space commutes with this flat extension, since each intersection
is a kernel. Let \(O_0\) be the original coefficient DVR. Completion
and extension to \(O\) are faithfully flat, and each length-one
factor gains the relative ramification index \(e\). Thus, for the
finite quotient \(Q\) in \eqref{eq:tag-c5},
\[
 \frac{\operatorname{length}_O(Q\otimes_{O_0}O)}
      {\operatorname{length}_O(O/pO)}
 =\frac{e\,\operatorname{length}_{O_0}Q}
       {e\,\operatorname{length}_{O_0}(O_0/pO_0)} .
\]
This gives the normalized length in the original degree; faithful
flatness also descends the inclusion.

\proofstep{Sum the excess on the selected weights.}

The excess \(2S_*\sigma\) in \eqref{eq:tag-c5} has the small total cost used
in Section~\ref{subsec:frobenius-selection}. The adjusted degrees at the tops satisfy
\[
  0\le a_s-2c_G\le6a/5.
\]
 The up gain is their
sum times \(p/(p^2-1)\) minus a sum of nonnegative minima and a
nonnegative logarithmic error. The price at a smaller prime is a
nonnegative multiple of a sum of the width and height bounds.

At
a larger prime its variable terms are the top filling minus the
base second weight, nonnegative by the chosen offsets and the
upper auxiliary allowance in the filling. Thus each term in \(S_v\)
before rounding is at most \(8N_va/(5p)\), by
\(p/(p^2-1)\le4/(3p)\). For \(N_v=0\) the ceiling is zero.
Otherwise 
\[
  a/p\ge b'\to\infty
\]
 on \(p<X'\), so rounding the
sum still gives \(S_v\le2mN_va/p\). With the selected weights,
\[
    \sum_{p\ {\rm served}}(2 S_*\sigma)\log p
             \le m\,|T|\,a\, x^{-9}.
                                                               \tag{C6}\label{eq:tag-c6}
\]
The set \(T\) here is the set of top positions. The estimate applies
on every subset of the primes and does not include a factor for the
number of words or coefficient places.

\proofstep{Bound the full baseline for the change of rank.}

It is useful to keep the larger term \(\tau\) in both the index
and the divisibility in \eqref{eq:tag-c5}. The crude bound

\[
  2\sum_p\tau(p)\log p\le b^2
\]
 suffices for the change in rank.
Indeed only \(p<X'=\exp(x^{\epsilon/3})\) can enter. Both order
formulas are at most \(x^{O(1)}a\) uniformly there. For the large
primes the sum of the truncated endpoint minima is at most \(2Na\);
the other terms consist of polynomial factors times \(a/p^2\),
\(\Lambda a/(pk)\) or \(Q(c_6+k)\), and a polynomial times \(c_G\).

Here \(Q^4(c_6+k)=o(b')\) up to polynomial factors by the
large-prime comparison in Section~\ref{sec:host-edge-tests},
since \(\log Q=o(x^\epsilon)\) and
\(\log(1+c_6)=o(x^\epsilon)\). The smaller-prime orders and both
depths obey the same crude bound by that section. The number of
orbits, \(m,f,k\) and \(\log a\) grow polynomially in \(x\). Thus
the sum is at most 
\[
  x^{O(1)}a X'\log(2X')=o(b^2),
\]
 since
\(a=xb\) and \(\log b=x^\epsilon\).

\end{proof}

\subsubsection{Mean and maximal slopes of the selected column}

The local indices give the change of determinant at all finite places. Combining them with the fixed shift compares the slopes of the resulting sublattice and its ambient column.

\begin{proposition}[Slope comparison on the shifted trace-zero column]\label{prop:shifted-column-slopes}
On the chosen word and its common division row, the local
conditions give the full unshifted sublattice \(N^{(0)}\) of
\(S_R\). The shift by \(b_w=\lfloor e_ws_u(w)\rfloor\) gives
\(N_u\) with the mean and maximal slope comparisons
\textup{\eqref{eq:tag-c7}} and \textup{\eqref{eq:tag-c8}} in the same induced metric.
Writing \(E_{\rm fl}\) for the floor correction in the common shift, its absolute size satisfies
\[
 |E_{\rm fl}|\le2(\log4+5+C_\Gamma)X_*.
\]

\end{proposition}

\begin{proof}
\proofstep{Assemble the full column from its completed lattices.}

Return to \(M_{\rm ram}=S_R\) over \(K\). At each place away from
\(T_u\) prescribe the lattice just computed, and use the image of
the product of the normalized up operators over \(T_u\) from
Proposition~\ref{prop:ramified-selected-hecke-colength}. These are sublattices of
the localization of \(S_R\), equal to it outside finitely many
places. They define a full global sublattice \(N^{(0)}\): take the
kernel of the map from \(S_R\) to the sum of all these finite
local quotients. Its localization is the prescribed kernel, since
localizing kills the quotients at the other primes and gives a
surjection at the given prime. A nonzero product of the local
annihilator ideals times \(S_R\) lies in the kernel.
Equivalently, if \(Q_w\) is the prescribed quotient at \(w\), localization gives
\[
  (S_R/N^{(0)})_w\simeq Q_w,\qquad \operatorname{rank}N^{(0)}=\operatorname{rank}S_R .
\]
 This also
proves full rank when some of the local kernels have positive
colength.

Let \(n_w\) be the length of the quotient in this
localization and \(d_R=d_{\rm ram}\). Set
\(\ell_w=\log {\rm N}w/[K:\mathbb Q]\), \(e_w=v_w(p)\). Denote the
total period valuation allowance in \(p\)-normalized units by
\(s_u(w)\). At a used place this is
\[
 s_u(w)=\sum_{j=1}^m
 \bigl(\mathrm{gain}^{\mathrm{keep}}_j(w)+F_j(w)-E_{u,j}(w)\bigr),
\]
where the first term is zero at a dropped prime. Use \(s_u(w)=0\)
off the finite support of these data. Its integral shift is

\[
  b_w=\lfloor e_ws_u(w)\rfloor.
\]

\proofstep{Identify the local tests in the period realization.}

The inclusion from the selected column space into \(M_u\)
has image the trace-zero summand. Its metric is the
restriction of \eqref{eq:selected-column-gram}. This follows
by evaluating along the same word inclusion both in the
complex period and in the Hermitian column form. In that
inclusion the initial integral lattice is the preimage
of \(M_u\), so it is \(S_R\) with the stated metric.

The prescribed local tests for \(\mathrm{NuLat}(u)\) are
the normalized up-operator image at \(T_u\), the orbit
image in the integral frame at a prime with a selected orbit, and the
initial integral condition at an prime without a selected orbit. They
are tests on the image under the same inclusion.
These give equality with \(N^{(0)}\) by localization.
For full lattices in the same rational space the local tests suffice by
\[
  L=\bigcap_{w\ {\rm finite}} L_w .
\]
In particular the evidence for the equality is the local
membership conditions, in addition to the quotient lengths.

The division package uses the pullback along its word-coordinate
isomorphism \(e_u\). The period identity on this isomorphism
uses its same assigned tuple and row evaluation on both
sides. The shifted column in the family is
\(\bigl(\prod_w\mathfrak p_w^{-b_w}\bigr)\mathrm{NuLat}(u)\).
Because \(e_u\) is \(K\)-linear, multiplication by this
fractional ideal commutes with the pullback. Thus the
column on which the period is tested is the shift of
\(N^{(0)}\), with the same metrics as the global slope.
\label{cor:column-orbit-length-1}

\proofstep{Compute the determinant degree and bound the maximal slope.}

The shifted lattice \(N_u\) of Section~\ref{subsec:division-shift}
then has the two comparisons
\[
 \begin{split}
 \mu(N_u)&=\mu(S_R)+2\sum_w\ell_w b_w-
                           \frac{2}{d_R}\sum_w\ell_w n_w,\\
 \mu_{\max}(N_u)&\le\mu_{\max}(S_R)
                     +2\sum_w\ell_w\bigl(b_w-e_w\tau(p)\bigr).
 \end{split}                                                     \tag{C7}\label{eq:tag-c7}
\]
Use \(\tau(T_u)=0\). The first line follows by the index on determinant
lattices and the degree of the shift ideal. For the second,
\eqref{eq:tag-c5} gives an inclusion of \(N_u\) into the ideal multiple
\(\prod_w\mathfrak p_w^{e_w\tau(p)-b_w}S_R\).
Every sublattice has its induced metric there; tensoring by an
invertible ideal adds the same constant to every slope.

\proofstep{Keep the full baseline in the common shift.}

To see the size of the difference, put \(B=\sum_w2\ell_we_w\tau(p)\)
and \(I_0=\sum_w2\ell_we_w(2S_*\sigma)\), setting the summands to zero
over \(T_u\). Let \(B_T=2\sum_{w\mid T_u}\ell_wn_w/d_R\).
With \(H=2\sum_w\ell_w b_w-B\), \eqref{eq:tag-c5} and \eqref{eq:tag-c7} give
\[
 \begin{split}
 \mu(N_u)&\ge\mu(S)+H-
  \left\{ \frac{f\log(T_u-1)}{T_u-2}
       +(d/d_R-1)B+(d/d_R)I_0+B_T\right\},\\
 \mu_{\max}(N_u)&\le\mu_{\max}(S)+H.
 \end{split}                                                     \tag{C8}\label{eq:tag-c8}
\]
The common term includes the subtraction of the \emph{full} baseline
index. Here \(\sum_{w\mid p}\ell_w e_w=\log p\) gives \(B\le b^2\)
and \(I_0\le2m|T|a x^{-9}\) by \eqref{eq:tag-c6}. The exact product rank gives

\[
  d/d_R\le 1+4f/T_u
\]
 for \(2f/T_u\le1/2\). Since \(T_u\ge b^{m/2}\),
\(m\ge6\) at late points, and \(f\) grows polynomially in \(x\), the
first two errors in braces together are eventually at most one.
The bound at the assigned prime gives
\(B_T\ll m |T|(a/k+b')\log T_u\) for the prescribed top weight
widths. All these bounds persist over the growing coefficient fields.

\proofstep{Bound the finite-support rounding correction.}

Above \(X_*=X''(n)\) the gain vanishes off the exceptional
prime set of the division estimate. There the forced term
is supported on \(g_0,g_1,d\).
Writing \(F_{\rm forced}(p)\) for this forced term, its support beyond the cutoff satisfies
\[
  p>X_*,\quad F_{\rm forced}(p)\ne0\quad\Longrightarrow\quad p\mid g_0g_1d.
\]
 The remaining exceptional
sets are contained in these primes, the prime assigned to the sequence and the
prime divisors of \(D_{\Gamma,u}\). Thus replacing
\(2\sum_w\ell_w b_w\) by \(2\sum_w\ell_we_ws_u(w)\) has absolute
error at most
\[
 2\big(X_*\log4+\log(g_0g_1dT_uD_{\Gamma,u})\big).
\]
We used \(\sum_{w\mid p}\ell_w\le\log p\), because \(e_w\ge1\),
and \(\prod_{p\le y}p\le4^y\), which follows by induction from the
middle binomial coefficients at \(2j\) and \(2j+1\).

The word size gives \(\log T_u\le2X_*\) and the denominator
estimate gives \(\log D_{\Gamma,u}\le C_\Gamma X_*\)
uniformly in the word. Each of \(\log g_0,\log g_1,\log d\)
is at most \(X_*\).
The bound for the three auxiliary logarithms is
\[
  \log(g_0g_1d)=\log g_0+\log g_1+\log d\le 3X_* .
\]
 Thus the absolute rounding error is at
most \(2(\log4+5+C_\Gamma)X_*\). The shift and this rounding
bound give the mean-slope comparison between \(N^{(0)}\)
and \(N_u\). The maximal slope in \eqref{eq:tag-c8} uses the full
baseline index before the bounded loss in the mean.
The valuation of the period on the generators of this
column gives its map to the dual of the row lattice with the prescribed local conditions by
Sections~\ref{subsec:hecke-period} and~\ref{subsec:division-shift}.
\label{cor:column-orbit-length-2}
\end{proof}
